\documentclass[11pt]{article}
\usepackage[margin=1in]{geometry}
\usepackage[T1]{fontenc}
\usepackage{lmodern}
\usepackage{microtype}
\usepackage{amsmath,amssymb}
\usepackage[dvipsnames]{xcolor}
\usepackage{enumitem}
\usepackage[colorlinks=true,linkcolor=MidnightBlue,urlcolor=MidnightBlue,citecolor=MidnightBlue]{hyperref}
\setlist[itemize]{leftmargin=1.4em,itemsep=0.25em,topsep=0.35em}
\setlength{\parindent}{0pt}
\setlength{\parskip}{0.62em}

\newcommand{\statusbox}[1]{%
  \begin{center}
  \fcolorbox{MidnightBlue}{MidnightBlue!6}{%
    \parbox{0.91\linewidth}{\centering\bfseries\color{MidnightBlue}#1}}
  \end{center}}

\title{Literature Status Note\\[0.25em]
\large The Lipschitz-domain oscillation LIL from arXiv:1203.5290}
\author{Run \texttt{fa\_banach\_001} --- lane 2}
\date{27 August 2026}

\begin{document}
\maketitle

\statusbox{LITERATURE ALREADY ANSWERED --- LIPSCHITZ-DOMAIN BRANCH ONLY}

\section*{Source question}

K. S. Eikrem, E. Malinnikova, and P. A. Mozolyako study harmonic functions
on the upper half-space with growth controlled by a decreasing doubling weight
$v$ and prove a law of the iterated logarithm for
\[
 I_u(x,s)=\int_s^1 u(x,t)\,d\!\left(\frac1{v(t)}\right).
\]
In Section 5, ``Concluding remarks and open problems,'' PDF page 19 of
arXiv:1203.5290, they ask whether their Theorem 5 extends to Lipschitz domains
\cite{EMM}.  In the same paragraph they state a separate, stronger local
conjecture under the radial hypothesis $|u(x,t)|\leq C v(t)$ for $x$ in a
positive-measure set.

\section*{Supporting answer}

P. Mozolyako's later paper explicitly identifies the source theorem and says:
``Our main goal is to extend this result to the domains ... above the graphs of
Lipschitz functions'' \cite[p.~2]{Moz}.  Its Theorem 1, on supporting PDF page
2, proves exactly that extension.  If
\[
 \Omega_\phi=\{(x,y):y\geq\phi(x)\},
\]
where $\phi$ is Lipschitz, and $u$ satisfies
\[
 |u(x,y)|\leq K w\!\left(\operatorname{dist}((x,y),\partial\Omega_\phi)\right),
\]
then for
\[
 I(x,\delta)=\int_\delta^1
 u(x,\phi(x)+y)\,d\!\left(\frac1{w(y)}\right)
\]
one has
\[
 \limsup_{\delta\to0}
 \frac{I(x,\delta)}
 {\sqrt{\log w(\delta)\,\log\log\log w(\delta)}}
 \leq C\,\|u\|_{w,\infty}
 \quad\text{for almost every }x\in\mathbb R^n.
\]
The supporting authors knew they were answering the source question: the paper
cites the source as [EMM], restates its upper-half-space theorem as Theorem A,
and presents Theorem 1 as the requested extension.  It also notes that the same
conclusion holds for certain star-like Lipschitz domains.

\section*{Scope boundary}

This identification does \emph{not} close the source paragraph's radial-only
local conjecture.  Theorem 2 of arXiv:1507.07455 proves a local result under the
stronger non-tangential hypothesis
\[
 |u(x,y)|\leq K w(y)\qquad\text{whenever }|x-x_0|\leq My,
 \quad x_0\in\Sigma.
\]
Immediately afterward, on supporting PDF page 3, the author states that it is
unknown whether this can be replaced by radial growth control.  Therefore the
durable classification is:

\begin{itemize}
  \item Lipschitz-domain extension: fully answered in later literature.
  \item Radial-only local version on a positive-measure set: still open in the
        decisive supporting paper; no claim of resolution is made here.
\end{itemize}

\section*{Evidence and review recommendation}

The identification used the arXiv source files and local PDFs for both papers,
plus exact-phrase and keyword searches for the source question.  The decisive
cross-checks are source PDF page 19, supporting PDF page 2 (Theorem 1), and
supporting PDF page 3 (the radial/non-tangential distinction).

Human review recommendation: classify only the Lipschitz-domain branch as an
already-answered literature result; retain the radial local conjecture as
unresolved unless a later exact source is found.

\begin{thebibliography}{9}
\bibitem{EMM}
K. S. Eikrem, E. Malinnikova, and P. A. Mozolyako,
\emph{Wavelet decomposition of harmonic functions in growth spaces},
arXiv:1203.5290; published as \emph{Wavelet characterization of growth spaces
of harmonic functions}, J. Anal. Math. 122 (2014), 87--111.

\bibitem{Moz}
P. Mozolyako,
\emph{Boundary oscillations of harmonic functions in Lipschitz domains},
arXiv:1507.07455 (2015).
\end{thebibliography}

\end{document}
